% figures/w02-pid-dialog.svg — Simulink PID Controller 블록에서 안티와인드업을 고르는 자리.
%
%   bash _tools/tikz2svg.sh figures/src/w02-pid-dialog.tex figures/w02-pid-dialog.svg
%
% 이것은 **스크린샷이 아니다.** CLAUDE.md 3 「GUI 창은 캡처할 수 없으므로 TikZ
% 개략도로 그리고 "스크린샷" 이라고 부르지 않는다」에 따른 개략도다. 필드 이름과
% 선택지 문자열은 R2023b~R2025a 의 PID Controller 블록에서 그대로 옮겼고,
% 창의 픽셀 배치를 흉내 내려 하지 않았다 — 학생이 찾아야 하는 것은 픽셀이 아니라
% **탭 이름과 필드 이름**이다.
%
% (a) back-calculation 을 골랐을 때  (b) clamping 을 골랐을 때.
% 둘의 차이는 Kb 칸이 나타나는가 하나뿐이고, 그것이 이 그림의 요점이다.
\documentclass[border=5pt]{standalone}
% gnc-style.tex — 이 볼트의 TikZ 그림이 공유하는 스타일.
%
% 저널 그림 규약을 스타일 하나로 굳혀 둔다. 그림마다 선 굵기나 화살촉을
% 다시 정하지 않는다 — 그것이 그림이 제각각으로 보이는 원인이다.
%
% 쓰는 법:  % gnc-style.tex — 이 볼트의 TikZ 그림이 공유하는 스타일.
%
% 저널 그림 규약을 스타일 하나로 굳혀 둔다. 그림마다 선 굵기나 화살촉을
% 다시 정하지 않는다 — 그것이 그림이 제각각으로 보이는 원인이다.
%
% 쓰는 법:  % gnc-style.tex — 이 볼트의 TikZ 그림이 공유하는 스타일.
%
% 저널 그림 규약을 스타일 하나로 굳혀 둔다. 그림마다 선 굵기나 화살촉을
% 다시 정하지 않는다 — 그것이 그림이 제각각으로 보이는 원인이다.
%
% 쓰는 법:  \input{gnc-style}  (figures/src/ 안에서)

\usepackage{tikz}
\usepackage{amsmath}
\usepackage{xcolor}
\usetikzlibrary{arrows.meta, positioning, calc, fit, backgrounds, decorations.pathmorphing}

% ---- 색 --------------------------------------------------------------------
% 색은 강조 하나에만 쓴다. 나머지는 검정.
\definecolor{ink}{HTML}{1A1A1A}
\definecolor{accent}{HTML}{7C3AED}
\definecolor{warn}{HTML}{B91C1C}
\definecolor{soft}{HTML}{666666}

% ---- 선과 화살촉 -----------------------------------------------------------
\tikzset{
  % 신호선. 화살촉은 작고 뾰족한 것 하나뿐이다.
  sig/.style   = {ink, line width=0.5pt, -{Stealth[length=4pt, width=3pt]}},
  wire/.style  = {ink, line width=0.5pt},
  % 강조 경로 — 파선으로 구분한다. 흑백 인쇄에서도 살아야 하기 때문이다.
  asig/.style  = {accent, line width=0.5pt, densely dashed,
                  -{Stealth[length=4pt, width=3pt]}},
  awire/.style = {accent, line width=0.5pt, densely dashed},
  % 블록. 흰 바탕, 직각, 검은 테두리.
  blk/.style   = {draw=ink, line width=0.55pt, fill=white,
                  minimum width=17mm, minimum height=8mm, inner sep=2pt, font=\small},
  % 합산점
  sum/.style   = {draw=ink, line width=0.55pt, fill=white, circle,
                  inner sep=0pt, minimum size=4.6mm, font=\scriptsize},
  asum/.style  = {draw=accent, line width=0.55pt, fill=white, circle,
                  inner sep=0pt, minimum size=4.6mm, font=\scriptsize},
  % 분기점
  dot/.style   = {ink, fill=ink, circle, inner sep=0pt, minimum size=1.5mm},
  % 신호 이름 — 선 위에 작게
  sname/.style = {font=\small, inner sep=1.5pt},
  % 그림 안의 짧은 설명. 그림 안의 글은 최소로 둔다
  note/.style  = {font=\scriptsize, soft, inner sep=1.5pt},
  anote/.style = {font=\scriptsize, accent, inner sep=1.5pt},
  % 축
  axis/.style  = {ink, line width=0.5pt},
  tick/.style  = {font=\scriptsize, soft},
}

% 캡션 한 줄 — 그림 아래 가는 선과 함께
\newcommand{\gnccaption}[2]{%
  \node[anchor=north west, text width=#1, align=left, font=\scriptsize, ink]
        at (current bounding box.south west) {\vspace{2pt}#2};}
  (figures/src/ 안에서)

\usepackage{tikz}
\usepackage{amsmath}
\usepackage{xcolor}
\usetikzlibrary{arrows.meta, positioning, calc, fit, backgrounds, decorations.pathmorphing}

% ---- 색 --------------------------------------------------------------------
% 색은 강조 하나에만 쓴다. 나머지는 검정.
\definecolor{ink}{HTML}{1A1A1A}
\definecolor{accent}{HTML}{7C3AED}
\definecolor{warn}{HTML}{B91C1C}
\definecolor{soft}{HTML}{666666}

% ---- 선과 화살촉 -----------------------------------------------------------
\tikzset{
  % 신호선. 화살촉은 작고 뾰족한 것 하나뿐이다.
  sig/.style   = {ink, line width=0.5pt, -{Stealth[length=4pt, width=3pt]}},
  wire/.style  = {ink, line width=0.5pt},
  % 강조 경로 — 파선으로 구분한다. 흑백 인쇄에서도 살아야 하기 때문이다.
  asig/.style  = {accent, line width=0.5pt, densely dashed,
                  -{Stealth[length=4pt, width=3pt]}},
  awire/.style = {accent, line width=0.5pt, densely dashed},
  % 블록. 흰 바탕, 직각, 검은 테두리.
  blk/.style   = {draw=ink, line width=0.55pt, fill=white,
                  minimum width=17mm, minimum height=8mm, inner sep=2pt, font=\small},
  % 합산점
  sum/.style   = {draw=ink, line width=0.55pt, fill=white, circle,
                  inner sep=0pt, minimum size=4.6mm, font=\scriptsize},
  asum/.style  = {draw=accent, line width=0.55pt, fill=white, circle,
                  inner sep=0pt, minimum size=4.6mm, font=\scriptsize},
  % 분기점
  dot/.style   = {ink, fill=ink, circle, inner sep=0pt, minimum size=1.5mm},
  % 신호 이름 — 선 위에 작게
  sname/.style = {font=\small, inner sep=1.5pt},
  % 그림 안의 짧은 설명. 그림 안의 글은 최소로 둔다
  note/.style  = {font=\scriptsize, soft, inner sep=1.5pt},
  anote/.style = {font=\scriptsize, accent, inner sep=1.5pt},
  % 축
  axis/.style  = {ink, line width=0.5pt},
  tick/.style  = {font=\scriptsize, soft},
}

% 캡션 한 줄 — 그림 아래 가는 선과 함께
\newcommand{\gnccaption}[2]{%
  \node[anchor=north west, text width=#1, align=left, font=\scriptsize, ink]
        at (current bounding box.south west) {\vspace{2pt}#2};}
  (figures/src/ 안에서)

\usepackage{tikz}
\usepackage{amsmath}
\usepackage{xcolor}
\usetikzlibrary{arrows.meta, positioning, calc, fit, backgrounds, decorations.pathmorphing}

% ---- 색 --------------------------------------------------------------------
% 색은 강조 하나에만 쓴다. 나머지는 검정.
\definecolor{ink}{HTML}{1A1A1A}
\definecolor{accent}{HTML}{7C3AED}
\definecolor{warn}{HTML}{B91C1C}
\definecolor{soft}{HTML}{666666}

% ---- 선과 화살촉 -----------------------------------------------------------
\tikzset{
  % 신호선. 화살촉은 작고 뾰족한 것 하나뿐이다.
  sig/.style   = {ink, line width=0.5pt, -{Stealth[length=4pt, width=3pt]}},
  wire/.style  = {ink, line width=0.5pt},
  % 강조 경로 — 파선으로 구분한다. 흑백 인쇄에서도 살아야 하기 때문이다.
  asig/.style  = {accent, line width=0.5pt, densely dashed,
                  -{Stealth[length=4pt, width=3pt]}},
  awire/.style = {accent, line width=0.5pt, densely dashed},
  % 블록. 흰 바탕, 직각, 검은 테두리.
  blk/.style   = {draw=ink, line width=0.55pt, fill=white,
                  minimum width=17mm, minimum height=8mm, inner sep=2pt, font=\small},
  % 합산점
  sum/.style   = {draw=ink, line width=0.55pt, fill=white, circle,
                  inner sep=0pt, minimum size=4.6mm, font=\scriptsize},
  asum/.style  = {draw=accent, line width=0.55pt, fill=white, circle,
                  inner sep=0pt, minimum size=4.6mm, font=\scriptsize},
  % 분기점
  dot/.style   = {ink, fill=ink, circle, inner sep=0pt, minimum size=1.5mm},
  % 신호 이름 — 선 위에 작게
  sname/.style = {font=\small, inner sep=1.5pt},
  % 그림 안의 짧은 설명. 그림 안의 글은 최소로 둔다
  note/.style  = {font=\scriptsize, soft, inner sep=1.5pt},
  anote/.style = {font=\scriptsize, accent, inner sep=1.5pt},
  % 축
  axis/.style  = {ink, line width=0.5pt},
  tick/.style  = {font=\scriptsize, soft},
}

% 캡션 한 줄 — 그림 아래 가는 선과 함께
\newcommand{\gnccaption}[2]{%
  \node[anchor=north west, text width=#1, align=left, font=\scriptsize, ink]
        at (current bounding box.south west) {\vspace{2pt}#2};}

% 체크박스 [x] 와 드롭다운 화살표에 amssymb 가 필요하다. gnc-style 은 amsmath 만
% 싣는다 — 공용 스타일에 넣지 않고 이 그림에서만 부른다.
\usepackage{amssymb}

\tikzset{
  win/.style   = {draw=ink, line width=0.6pt, fill=white, rounded corners=0.6mm},
  bar/.style   = {draw=ink, line width=0.5pt, fill=black!6},
  fld/.style   = {draw=soft, line width=0.4pt, fill=white,
                  minimum height=4.2mm, inner xsep=1.2mm, font=\scriptsize, anchor=west},
  afld/.style  = {draw=accent, line width=0.7pt, fill=accent!6,
                  minimum height=4.2mm, inner xsep=1.2mm, font=\scriptsize,
                  anchor=west, text=accent},
  lbl/.style   = {font=\scriptsize, ink, anchor=east},
  tab/.style   = {font=\scriptsize, ink, anchor=west},
  atab/.style  = {font=\scriptsize\bfseries, accent, anchor=west},
}

\begin{document}
\begin{tikzpicture}[x=1mm, y=1mm]

\node[font=\small\bfseries, anchor=north west] at (0,10)
     {Where the two schemes are selected in the PID Controller block};
\node[anchor=north west, align=left, text width=178mm, font=\scriptsize, soft] at (0,4) {%
  A schematic of the dialog, not a screenshot. Both schemes live on the same tab,
  chosen from one menu; picking \emph{back-calculation} is what makes the
  $K_b$ field appear. Everything else on the tab is identical.};

% =====================================================================
\def\WW{86}      % 창 너비
\def\HH{50}      % 창 높이

% ---------- (a) ------------------------------------------------------
\begin{scope}[shift={(0,-12)}]
  \node[font=\scriptsize\bfseries, anchor=west, text=ink] at (0,2)
       {(a)\quad \texttt{aw\_mode = back-calculation}};

  \draw[win] (0,0) rectangle (\WW,-\HH);
  \draw[bar] (0,0) rectangle (\WW,-6);
  \node[tab, font=\scriptsize\bfseries] at (2,-3) {Block Parameters: PID Controller};

  % 탭 줄
  \draw[soft, line width=0.4pt] (0,-11) -- (\WW,-11);
  \node[tab]  at (2,-8.5)  {Main};
  \node[atab] at (18,-8.5) {PID Advanced};
  \node[tab]  at (46,-8.5) {Data Types};
  \draw[accent, line width=1.0pt] (17,-11) -- (43,-11);

  \node[lbl] at (40,-16) {Output saturation:};
  \node[fld, minimum width=40mm] at (42,-16) {Limit output\quad$\boxtimes$};

  \node[lbl] at (40,-22) {Upper limit:};
  \node[fld, minimum width=40mm] at (42,-22) {\texttt{X\_max}};

  \node[lbl] at (40,-28) {Lower limit:};
  \node[fld, minimum width=40mm] at (42,-28) {\texttt{X\_min}};

  \node[lbl] at (40,-35) {Anti-windup method:};
  \node[afld, minimum width=40mm] at (42,-35) {back-calculation\hfill$\blacktriangledown$};

  \node[lbl] at (40,-41) {Back-calc. coefficient (Kb):};
  \node[afld, minimum width=40mm] at (42,-41) {\texttt{1/tau\_u}};

  \node[anote, anchor=north west, text width=40mm, align=left] at (42,-44.2)
       {appears \emph{only} for back-calculation};
\end{scope}

% ---------- (b) ------------------------------------------------------
\begin{scope}[shift={(96,-12)}]
  \node[font=\scriptsize\bfseries, anchor=west, text=ink] at (0,2)
       {(b)\quad \texttt{aw\_mode = clamping}};

  \draw[win] (0,0) rectangle (\WW,-\HH);
  \draw[bar] (0,0) rectangle (\WW,-6);
  \node[tab, font=\scriptsize\bfseries] at (2,-3) {Block Parameters: PID Controller};

  \draw[soft, line width=0.4pt] (0,-11) -- (\WW,-11);
  \node[tab]  at (2,-8.5)  {Main};
  \node[atab] at (18,-8.5) {PID Advanced};
  \node[tab]  at (46,-8.5) {Data Types};
  \draw[accent, line width=1.0pt] (17,-11) -- (43,-11);

  \node[lbl] at (40,-16) {Output saturation:};
  \node[fld, minimum width=40mm] at (42,-16) {Limit output\quad$\boxtimes$};

  \node[lbl] at (40,-22) {Upper limit:};
  \node[fld, minimum width=40mm] at (42,-22) {\texttt{X\_max}};

  \node[lbl] at (40,-28) {Lower limit:};
  \node[fld, minimum width=40mm] at (42,-28) {\texttt{X\_min}};

  \node[lbl] at (40,-35) {Anti-windup method:};
  \node[afld, minimum width=40mm] at (42,-35) {clamping\hfill$\blacktriangledown$};

  %  (a) 에서 Kb 칸이 있던 자리. 비어 있다는 것 자체가 이 판의 요점이다.
  \draw[soft, line width=0.4pt, densely dashed] (42,-41) rectangle (82,-38.9);
  \node[note, anchor=north west, text width=40mm, align=left] at (42,-44.2)
       {no coefficient to set — clamping has nothing to tune};
\end{scope}

\end{tikzpicture}
\end{document}
